\documentclass[11pt]{article}
\usepackage{amsmath,amssymb,amsthm}

\newtheorem{theorem}{Theorem}
\newtheorem{lemma}[theorem]{Lemma}
\newtheorem{corollary}[theorem]{Corollary}

\newcommand{\Up}{U_+}

\title{Transport of the deficit-four horizon along fixed-seed returns}
\author{}
\date{27 August 2026}

\begin{document}
\maketitle

Let
\[
  x_{t+1}=\Up(x_t):=\frac{3x_t+1}{2^{a_t}},
  \qquad
  a_t:=v_2(3x_t+1),
  \qquad
  A_t:=\sum_{i=0}^{t-1}a_i,
  \tag{1}
\]
where every $x_t$ is a positive odd integer.  We assume that every prefix
from $x_0$ is a strict barrier:
\[
  2^{A_t}<3^t
  \qquad(t\geq1).
  \tag{2}
\]

\begin{lemma}[exact return lift and extension law]
Suppose that the length-$L$ exponent prefix occurs again at position $d>0$:
\[
  (a_d,\ldots,a_{d+L-1})=(a_0,\ldots,a_{L-1}).
  \tag{3}
\]
Then there is an integer $t_L\geq1$ such that
\[
\boxed{
  x_d=x_0+t_L2^{A_L+1},
  \qquad
  x_{d+L}=x_L+2t_L3^L.}
\tag{4}
\]
For $h\geq0$, put
\[
  B_h:=A_{L+h}-A_L.
  \tag{5}
\]
The occurrence in (3) extends through the next $h$ exponent letters if and
only if
\[
\boxed{2^{B_h}\mid t_L.}
\tag{6}
\]
Whenever (6) holds,
\[
\boxed{
  x_{d+L+h}=x_{L+h}
  +2\frac{t_L}{2^{B_h}}3^{L+h}.}
\tag{7}
\]
In particular,
\[
  \min\{v_3(x_{d+L+h}+1),L+h\}
  =\min\{v_3(x_{L+h}+1),L+h\}.
  \tag{8}
\]
\end{lemma}

\begin{proof}
A positive exponent word of total exponent $A$ determines one odd seed class
modulo $2^{A+1}$.  We include the standard induction.  Suppose a prefix of
length $k$ and total exponent $A$ has representative $\rho$.  Replacing its
initial value by $\rho+\ell2^{A+1}$ replaces its endpoint by
\[
  x_k(\rho)+2\ell3^k.
\]
Appending an exponent $a$ requires
\[
  3x_k(\rho)+1+6\ell3^k\equiv2^a\pmod {2^{a+1}}.
\]
After division by $2$, the coefficient of $\ell$ is odd, so this congruence
selects one class $\ell\pmod {2^a}$.  This proves the cylinder assertion.

The two starting states in (3) therefore agree modulo $2^{A_L+1}$, which
gives the first identity in (4).  Condition (2) and the positive affine carry
give $x_d>x_0$, so $t_L\geq1$.  Applying the common length-$L$ affine map to
the starting-state difference gives the second identity.

The occurrence extends through $h$ further letters exactly when $x_d$ and
$x_0$ lie in the same length-$(L+h)$ cylinder.  By the first identity in
(4), this is equivalent to
\[
  2^{A_{L+h}+1}\mid t_L2^{A_L+1},
\]
which is (6).  Applying the common length-$(L+h)$ affine map proves (7), and
reduction of (7) modulo $3^{L+h}$ proves (8).
\end{proof}

Let
\[
  \mathcal S_4:=
  \{\varnothing,12,112,1112,1113,1122,1212,11113\}
  \tag{9}
\]
be the proved deficit-four stem set.  A stem $s$ matches $x_0$ if it is the
initial exponent word of $x_0$.  For a matching stem of length $h$, put
\[
  \nu_s:=v_2(x_h+1).
  \tag{10}
\]
The proved fixed-start horizon is
\[
  R_4(x_0):=
  \max_{\substack{s\in\mathcal S_4\\s\text{ matches }x_0}}
  (|s|+\nu_s-1).
  \tag{11}
\]
Define also the finite transport threshold
\[
  \Gamma_4(x_0):=
  \max_{\substack{s\in\mathcal S_4\\s\text{ matches }x_0}}
  (A_{|s|}+\nu_s).
  \tag{12}
\]

\begin{theorem}[fixed-seed horizon transport]
Assume (3), $L\geq5$, and
\[
  A_L+1>\Gamma_4(x_0).
  \tag{13}
\]
Then
\[
\boxed{R_4(x_d)=R_4(x_0).}
\tag{14}
\]
Since $A_L\geq L$, condition (13) holds for every sufficiently long returned
prefix of the fixed seed $x_0$.
\end{theorem}

\begin{proof}
Every stem in $\mathcal S_4$ has length at most $5$.  Thus (3) and $L\geq5$
show that exactly the same stems match $x_d$ and $x_0$.  Fix one such stem
$s$ and put $h:=|s|$.  Applying its common affine map to the first identity
in (4) gives
\[
  x_{d+h}-x_h=t_L3^h2^{A_L+1-A_h}.
  \tag{15}
\]
By (12) and (13),
\[
  A_L+1-A_h>\nu_s=v_2(x_h+1).
\]
Equation (15) therefore implies
\[
  v_2(x_{d+h}+1)=v_2(x_h+1)=\nu_s.
  \tag{16}
\]
Every term in the maximum defining $R_4$ is the same at the two starts, so
(14) follows.
\end{proof}

\begin{corollary}[returned-start exclusion]
Under the hypotheses of Theorem~2, suppose also that
\[
  L>R_4(x_0).
  \tag{17}
\]
Let $r\geq L$.  If the length-$r$ exponent block beginning at $x_d$ is a
strict barrier block, then
\[
\boxed{
  v_3(x_{d+r}+1)\leq r-5.}
\tag{18}
\]
In particular, if $x_d$ begins an infinite strict-barrier ray, then no
length $r\geq L$ from this returned start can be a deficit-four
synchronization hit.  Any such hit obtained from the return construction
must change the block start.
\end{corollary}

\begin{proof}
Theorem~2 gives $R_4(x_d)=R_4(x_0)<L\leq r$.  The fixed-start deficit-four
horizon theorem at $x_d$ now gives (18).  The remaining assertions follow
immediately.
\end{proof}

\begin{theorem}[first-branch exponent law]
Assume the hypotheses of Theorem~2.  Let $h$ be a finite maximal common extension
of the returned prefix, and put
\[
  N:=L+h,
  \qquad
  B:=A_N-A_L,
  \qquad
  T:=\frac{t_L}{2^B}.
  \tag{19}
\]
Let
\[
  y:=x_N,
  \qquad
  y':=x_{d+N}=y+2T3^N,
  \tag{20}
\]
and denote their distinct next exponents by
\[
  a:=v_2(3y+1),
  \qquad
  a':=v_2(3y'+1).
  \tag{21}
\]
Put $c:=v_2(T)$ and $z:=(3y+1)/2^a$.  Then $c<a$, and the first branch is
governed exactly by
\[
\boxed{
  \begin{aligned}
  c+1<a&\quad\Longrightarrow\quad a'=c+1,\\
  c+1=a&\quad\Longrightarrow\quad
  a'=a+v_2\left(z+\frac{T}{2^{a-1}}3^{N+1}\right)>a.
  \end{aligned}}
\tag{22}
\]
\end{theorem}

\begin{proof}
Maximality of $h$ and the extension law (6) give
\[
  2^B\mid t_L,
  \qquad
  2^{B+a}\nmid t_L,
\]
so $c<a$.  Equation (20) gives
\[
  3y'+1=2^az+2T3^{N+1}.
  \tag{23}
\]
The two terms on the right have $2$-adic valuations $a$ and $c+1$.  If they
are distinct, their sum has the smaller valuation.  If they are equal,
factoring out $2^a$ gives the second line of (22); the two remaining
summands are odd, so the displayed valuation is positive.
\end{proof}

\end{document}
